\chapter{最优化}
\label{chap:optimization}

第~\ref{chap:mathematical-analysis}章说明了函数在一点附近怎样变化。本章要把这些局部信息
用于一个整体问题：在所有允许的参数中，怎样判断某个参数已经足够好，又怎样逐步找到它？
贯穿全章的对象是带正则项的最小二乘
\begin{equation}
  F(\bm{w})
  =
  \frac{1}{2n}\lVert \bm{X}\bm{w}-\bm{y}\rVert_2^2
  +\frac{\lambda}{2}\lVert\bm{w}\rVert_2^2,
  \qquad
  \bm{X}\in\mathbb{R}^{n\times d},\quad \lambda\geq 0,
  \label{eq:opt-regularized-least-squares}
\end{equation}
以及线性资源约束
\begin{equation}
  \bm{A}\bm{w}=\bm{b},\qquad
  \bm{G}\bm{w}\leq\bm{h}.
  \label{eq:opt-linear-resource-constraints}
\end{equation}
式~\eqref{eq:opt-regularized-least-squares}与第~\ref{chap:sets-functions-and-proofs}章的
式~\eqref{eq:prelim-regularized-least-squares}是同一个目标的矩阵写法。正则项限制参数
规模，约束则直接排除超出预算的候选；两者看似都在“限制模型”，却改变了不同的数学对象。

本章先区分解、最优性条件与求解算法，再讨论曲率如何决定可以得到何种保证。约束引出乘子、
对偶和可核验的证书，非光滑项与不完整梯度则要求改变更新方式。分块、多目标和双层问题不是
一条由旧到新的算法路线，而是三种不同的问题组织方式。读完本章，读者应能判断一个优化结论
依赖哪些条件，并能指出条件失效时究竟失去了存在性、全局性、唯一性还是收敛速度。

\section{问题设定与解的性质}
\label{sec:opt-problem-and-solutions}

一个优化问题必须同时给出目标函数、变量和可行域。记
\[
  \min_{\bm{w}\in C}F(\bm{w}),\qquad
  C=\{\bm{w}\in\mathbb{R}^d:
  \bm{A}\bm{w}=\bm{b},\ \bm{G}\bm{w}\leq\bm{h}\}.
\]
若$C=\varnothing$，问题\term{不可行}（infeasible）；若$C\neq\varnothing$但
$\inf_{\bm{w}\in C}F(\bm{w})=-\infty$，问题\term{无下界}（unbounded below）。
只有可行且有有限下界，才值得继续问下确界是否由某个点达到。

\term{全局最优解}（global minimizer）$\bm{w}^{\star}\in C$满足
$F(\bm{w}^{\star})\leq F(\bm{w})$对所有$\bm{w}\in C$成立。\term{局部最优解}
（local minimizer）只要求该不等式在$\bm{w}^{\star}$的某个邻域内成立。
无约束可微问题的\term{驻点}（stationary point）满足
$\nabla F(\bm{w})=\bm{0}$。驻点只是局部一阶变化消失，可能是极小点、极大点或
第~\ref{sec:analysis-second-order-curvature}节所说的微分鞍点。
给定$\varepsilon>0$，若
\[
  F(\widehat{\bm{w}})-\inf_{\bm{w}\in C}F(\bm{w})\leq\varepsilon,
\]
就称$\widehat{\bm{w}}$为目标值意义下的$\varepsilon$近似解。对无约束可微问题，
另一种常见标准是
$\lVert\nabla F(\widehat{\bm{w}})\rVert_2\leq\varepsilon$；在没有额外曲率条件时，
这两种近似不能互换。一般约束问题的精确最优点也可能有非零梯度，此时应改用投影梯度、
可行方向导数或KKT残差等同时反映可行性与驻点性的标准。

若$C\subseteq\mathbb{R}^d$非空且紧，$F:C\to\mathbb{R}$连续，则
定理~\ref{thm:analysis-extreme-value}保证$F$在$C$上达到全局最小值。这里直接复用
第~\ref{chap:mathematical-analysis}章的极值存在结论。

可行域不紧时，目标本身也可能阻止最小化序列逃向无穷远。若
$\lVert\bm{w}\rVert_2\to\infty$时$F(\bm{w})\to\infty$，称$F$是
\term{强制的}（coercive）。连续性并不是这里所需的最弱条件。\term{下半连续}
（lower semicontinuous）要求对每个收敛序列
$\bm{w}^{(k)}\to\bm{w}$都有
\[
  F(\bm{w})\leq\liminf_{k\to\infty}F(\bm{w}^{(k)}).
\]
它允许函数值在极限点向下跳，却不允许极限点的函数值突然高于邻近值。等价地，每个
下水平集$\{\bm{w}:F(\bm{w})\leq a\}$都是闭集。

设$C$非空且闭，$F:C\to(-\infty,+\infty]$下半连续、强制，并且不恒为
$+\infty$且从不取$-\infty$。取$\bar{\bm{w}}\in C$使
$F(\bar{\bm{w}})<+\infty$，再把最小化序列限制到下水平集
\[
  K=\{\bm{w}\in C:F(\bm{w})\leq F(\bar{\bm{w}})\}.
\]
若$\inf_C F=F(\bar{\bm{w}})$，则$\bar{\bm{w}}$已经是最小点；否则
$\inf_C F<F(\bar{\bm{w}})$，任一最小化序列从某项起都落入$K$，可以只保留这段尾序列。
强制性使$K$有界，下半连续性与$C$的闭性使$K$闭，故$K$在有限维空间中紧。该尾序列
有子列收敛到$\bm{w}^{\star}\in K$，而下半连续性给出
$F(\bm{w}^{\star})\leq\liminf_kF(\bm{w}^{(k)})=\inf_C F$。因此下确界达到。
这个证明也指出各条件的职责：强制性阻止序列逃向无穷远，闭性保留可行极限，下半连续性
把目标值比较传到极限点。
当$\lambda>0$时，式~\eqref{eq:opt-regularized-least-squares}满足
$F(\bm{w})\geq\lambda\lVert\bm{w}\rVert_2^2/2$，所以即使可行域无界也不会沿参数
无穷大的方向降低目标。

存在不等于唯一。若$\lambda=0$且$\bm{X}$有非零零空间向量$\bm{v}$，那么
$F(\bm{w}+t\bm{v})=F(\bm{w})$，最小点可能形成一条直线。即使唯一，数值问题也可能
非常敏感。例如令$0<\epsilon\ll1$并考虑
\[
  f_b(w_1,w_2)=\frac12(\epsilon w_1^2+w_2^2)-bw_1.
\]
它对每个$b$都有唯一最小点$(b/\epsilon,0)$，但$b$改变$\Delta b$会使第一坐标改变
$\Delta b/\epsilon$。更一般的二次目标具有如下曲率矩阵：
\[
  \bm{H}=\frac{1}{n}\bm{X}^{\top}\bm{X}+\lambda\bm{I}
\]
当$\bm{H}\succ\bm{0}$时，若线性项受到加性扰动$\delta\bm{g}$，相应解变化满足
\[
  \lVert\delta\bm{w}\rVert_2
  \leq
  \frac{1}{\lambda_{\min}(\bm{H})}
  \lVert\delta\bm{g}\rVert_2.
\]
因此$1/\lambda_{\min}(\bm{H})$控制固定尺度加性扰动的绝对放大；经过目标和变量尺度
归一化后，条件数
$\kappa(\bm{H})=\lambda_{\max}(\bm{H})/\lambda_{\min}(\bm{H})$
才刻画相对病态程度。例如$\bm{H}=\epsilon\bm{I}$的条件数为$1$，但绝对扰动仍会被
放大$1/\epsilon$倍。唯一性回答“解有几个”，敏感性还取决于所采用的扰动尺度，
不能只由唯一性或条件数笼统判断。

\section{凸集与凸函数}
\label{sec:opt-convexity}

\subsection{线段、曲率与全局下界}

\term{凸集}（convex set）包含任意两点之间的整条线段。也就是说，对
$\bm{u},\bm{v}\in C$和$\theta\in[0,1]$，
$\theta\bm{u}+(1-\theta)\bm{v}\in C$。线性等式、线性不等式和范数球的解集都是凸集，
它们的交集仍为凸集，所以式~\eqref{eq:opt-linear-resource-constraints}定义凸可行域。

定义在凸集$C$上的函数$f$若满足
\begin{equation}
  f(\theta\bm{u}+(1-\theta)\bm{v})
  \leq
  \theta f(\bm{u})+(1-\theta)f(\bm{v}),
  \label{eq:opt-convex-function}
\end{equation}
就称为\term{凸函数}（convex function）。它的函数图像位于任意弦段下方，因此局部低点
不会被更远处的谷底击败\scite{math-boyd2004convex}

\begin{theorem}[凸函数的一阶与二阶判据]
\label{thm:opt-convexity-criteria}
设$C$为开凸集。可微函数$f:C\to\mathbb{R}$凸，当且仅当对任意
$\bm{u},\bm{v}\in C$，
\begin{equation}
  f(\bm{v})\geq f(\bm{u})
  +\langle\nabla f(\bm{u}),\bm{v}-\bm{u}\rangle.
  \label{eq:opt-first-order-convexity}
\end{equation}
若$f$二阶可微，则这又等价于
$\nabla^2f(\bm{w})\succeq\bm{0}$对所有$\bm{w}\in C$成立。
\end{theorem}

\begin{proof}
若$f$凸，令$\phi(t)=f(\bm{u}+t(\bm{v}-\bm{u}))$。由凸性，
$[\phi(t)-\phi(0)]/t\leq\phi(1)-\phi(0)$。令$t\downarrow0$便得
式~\eqref{eq:opt-first-order-convexity}。反过来，令
$\bm{z}=\theta\bm{u}+(1-\theta)\bm{v}$，分别把一阶不等式用于
$(\bm{z},\bm{u})$和$(\bm{z},\bm{v})$，乘以$\theta$与$1-\theta$后相加，
线性项抵消，得到式~\eqref{eq:opt-convex-function}。

若$f$凸，把一阶判据用于$\bm{w}$与$\bm{w}+t\bm{v}$，再对$t$求二阶局部极限，
得到$\bm{v}^{\top}\nabla^2f(\bm{w})\bm{v}\geq0$。反之，沿任意线段定义
$\phi(t)$，则
$\phi''(t)=(\bm{v}-\bm{u})^{\top}\nabla^2f(\bm{u}+t(\bm{v}-\bm{u}))
(\bm{v}-\bm{u})\geq0$，一元函数$\phi$凸，故$f$凸。
\end{proof}

反复应用二点凸性可得\term{Jensen不等式}（Jensen's inequality）。若
$\alpha_i\geq0$且$\sum_i\alpha_i=1$，则
\begin{equation}
  f\left(\sum_{i=1}^m\alpha_i\bm{w}_i\right)
  \leq\sum_{i=1}^m\alpha_i f(\bm{w}_i).
  \label{eq:opt-jensen}
\end{equation}
归纳步骤把前$m-1$个点先归一化为一个加权平均，再与第$m$个点应用二点凸性。
这个证明也说明非负和归一化不是装饰条件；允许负权重后，组合点可能落到凸包之外，
因而不能再由凸性保证Jensen不等式。

凸性还把最优性变成可核验的一阶条件。若$f$可微且凸，则
$\nabla f(\bm{w}^{\star})=\bm{0}$代入式~\eqref{eq:opt-first-order-convexity}
便得$f(\bm{v})\geq f(\bm{w}^{\star})$。因此凸问题的驻点是全局最小点。

严格凸要求非平凡线段上的式~\eqref{eq:opt-convex-function}严格成立，它保证至多一个
最小点，却不给统一曲率下界。\term{$\mu$强凸}（$\mu$-strongly convex）要求
\[
  f(\bm{v})\geq f(\bm{u})+
  \langle\nabla f(\bm{u}),\bm{v}-\bm{u}\rangle
  +\frac{\mu}{2}\lVert\bm{v}-\bm{u}\rVert_2^2.
\]
\term{$L$光滑}（$L$-smooth）则要求梯度满足
$\lVert\nabla f(\bm{u})-\nabla f(\bm{v})\rVert_2
\leq L\lVert\bm{u}-\bm{v}\rVert_2$。强凸给下曲率，光滑给上曲率；$x^4$严格凸但
不在全实线上强凸，也不在全实线上具有有限的光滑常数。

\subsection{最近点、分离与支撑}

\begin{theorem}[闭凸集外一点的严格分离]
\label{thm:opt-projection-separation}
设$C\subseteq\mathbb{R}^d$非空、闭且凸，$\bm{z}\notin C$。则存在
$\bm{a}\neq\bm{0}$和$c\in\mathbb{R}$，使
\[
  \langle\bm{a},\bm{x}\rangle\leq c
  <\langle\bm{a},\bm{z}\rangle
  \quad\text{对所有 }\bm{x}\in C.
\]
\end{theorem}

\begin{proof}
先证明所需最近点的存在性与唯一性。任取$\bm{x}_0\in C$并令
$R=\lVert\bm{z}-\bm{x}_0\rVert_2$。集合
\[
  C_R=C\cap\{\bm{x}:\lVert\bm{z}-\bm{x}\rVert_2\leq R\}
\]
非空、闭且有界，因而在有限维空间中紧。连续函数
$\lVert\bm{z}-\bm{x}\rVert_2^2$在$C_R$上达到最小值；球外各点到$\bm{z}$的
距离大于$R$，而$\bm{x}_0$的距离等于$R$，所以该最小点也是$C$上的最近点。
若$\bm{p}_1,\bm{p}_2$都是最近点，最近距离为$\delta$，则凸性保证其中点属于$C$，且
\[
  \left\lVert\bm{z}-\frac{\bm{p}_1+\bm{p}_2}{2}\right\rVert_2^2
  =
  \frac12\lVert\bm{z}-\bm{p}_1\rVert_2^2
  +\frac12\lVert\bm{z}-\bm{p}_2\rVert_2^2
  -\frac14\lVert\bm{p}_1-\bm{p}_2\rVert_2^2.
\]
当$\bm{p}_1\neq\bm{p}_2$时，右端严格小于$\delta^2$，与最近距离的定义矛盾，故最近点
唯一，记为$\bm{p}$。

对任意$\bm{x}\in C$，
$\bm{p}+t(\bm{x}-\bm{p})\in C$。函数
$q(t)=\lVert\bm{z}-\bm{p}-t(\bm{x}-\bm{p})\rVert_2^2$在$t=0$
处右导数非负，所以
$\langle\bm{z}-\bm{p},\bm{x}-\bm{p}\rangle\leq0$。取
$\bm{a}=\bm{z}-\bm{p}$和$c=\langle\bm{a},\bm{p}\rangle$，则
$\langle\bm{a},\bm{x}\rangle\leq c$，而
$\langle\bm{a},\bm{z}\rangle=c+\lVert\bm{a}\rVert_2^2>c$。
\end{proof}

等号对应的超平面若接触凸集，称为\term{支撑超平面}（supporting hyperplane）。
具体地说，若$\bm{p}$位于闭凸集$C$的边界，可取一列集外点
$\bm{z}^{(k)}\to\bm{p}$，对每个点应用严格分离并把法向量归一化。单位球的紧性给出
收敛子列，极限法向量$\bm{a}\neq\bm{0}$满足
$\langle\bm{a},\bm{x}\rangle\leq\langle\bm{a},\bm{p}\rangle$对所有
$\bm{x}\in C$成立。这说明边界处至少有一个支撑方向，但边界有棱角时法向量未必唯一。
分离把“点不属于集合”转化为一个线性不等式证书，但它本身不自动给出一般凸问题的强对偶；
还需要闭性或约束资格条件。

对称矩阵集合中的\term{半正定锥}（positive semidefinite cone）
\[
  \mathbb{S}_+^d=\{\bm{M}=\bm{M}^{\top}:
  \bm{v}^{\top}\bm{M}\bm{v}\geq0,\ \forall\bm{v}\}
\]
是凸锥，因为非负线性组合仍保持二次型非负。它也是闭集：若
$\bm{M}^{(k)}\to\bm{M}$且每个$\bm{M}^{(k)}\succeq\bm{0}$，则对任意固定
$\bm{v}$都有
$\bm{v}^{\top}\bm{M}\bm{v}
=\lim_k\bm{v}^{\top}\bm{M}^{(k)}\bm{v}\geq0$。
线性规划的非负正交锥和半正定规划的矩阵约束直接定义凸可行域；标准凸二次规划中的
半正定曲率矩阵则保证目标凸。三者共享的是凸目标与凸可行域的语言。若要进一步统一成
带线性目标的锥规划，还需引入辅助变量，对二次目标作上图或锥提升。

\section{较弱曲率与非凸结构}
\label{sec:opt-weaker-curvature}

\subsection{弱凸与补足曲率}

强凸同时控制函数形状和误差。实际目标可能不凸，却仍有足够结构支持算法分析。
\term{$\rho$弱凸}（$\rho$-weakly convex）表示
$f(\bm{w})+\rho\lVert\bm{w}\rVert_2^2/2$是凸函数。若$f$二阶可微，这等价于
$\nabla^2f(\bm{w})\succeq-\rho\bm{I}$：负曲率可以存在，但有统一下界。
当$\rho>0$且$0<\eta<1/\rho$时，对任意中心$\bm{z}$，邻近子问题的目标
\[
  f(\bm{w})+\frac{1}{2\eta}\lVert\bm{w}-\bm{z}\rVert_2^2
\]
是$(1/\eta-\rho)$强凸函数；若$\rho=0$，任意$\eta>0$都具有同样的补足曲率作用。
这正是弱凸性与第~\ref{sec:opt-nonsmooth-proximal}节邻近方法的联系。补入二次曲率
只让这个子问题更易控制，并没有把原函数变成凸函数；弱凸目标仍可有局部极小点和鞍点，
也不保证求得原问题的全局最优解。

\subsection{梯度控制目标误差的PL条件}

\term{Polyak--Łojasiewicz条件}（PL condition）不约束函数图像必须凸，而要求
\begin{equation}
  \frac{1}{2}\lVert\nabla f(\bm{w})\rVert_2^2
  \geq \mu\bigl(f(\bm{w})-f^{\star}\bigr),
  \qquad \mu>0.
  \label{eq:opt-pl-condition}
\end{equation}
它直接把小梯度与小目标误差联系起来。强凸可推出PL条件：对$\mu$强凸可微函数，把强凸
下界右侧关于$\bm{v}$最小化，得到
\[
  f^\star\geq
  f(\bm{w})-\frac{1}{2\mu}\lVert\nabla f(\bm{w})\rVert_2^2,
\]
移项正是式~\eqref{eq:opt-pl-condition}。这个推导使用全局强凸下界；只有局部Hessian
正定时，不能直接得到全域PL条件。反过来，PL函数可以有多个最小点。

同一个低维族能显示这些性质各自控制什么。令
\[
  f_a(x_1,x_2)=\frac12(x_1+x_2)^2+\frac{a}{2}(x_1-x_2)^2.
\]
当$a>0$时函数强凸；$a=0$时它凸且满足PL条件，但整条直线$x_1+x_2=0$都是最小点；
$a<0$时函数弱凸却无下界。若限制在紧可行域上，$a<0$时仍有最小点，但驻点
$(0,0)$可能是鞍点。这个例子说明弱凸、PL和强凸不是一条简单的强弱链：弱凸限制最坏负
曲率，PL控制梯度与目标误差，强凸还控制参数距离与唯一性。

PL条件也确实可以出现在非凸函数中。考虑
\[
  f(x,y)=x^2+2\sin^2x.
\]
它与$y$无关，因而最小点不唯一。其一阶导数为
$\partial f/\partial x=2(x+\sin 2x)$；当$x>0$时，若$x\leq1$则两项均非负，
若$x>1$则$x+\sin2x\geq x-1>0$，负半轴由奇对称性处理，所以导数只在$x=0$时
为零。比值
\[
  R(x)=\frac{\lVert\nabla f(x,y)\rVert_2^2}{2f(x,y)}
\]
在$x\neq0$时为正，并且$x\to0$时趋于$6$，$|x|\to\infty$时趋于$2$。
把两端之外的剩余区域限制为紧集，连续函数$R$在其中有正下界，故某个$\mu>0$使
式~\eqref{eq:opt-pl-condition}成立。另一方面，
$\partial^2 f/\partial x^2=2+4\cos2x$在某些点为负，所以$f$并不凸；该二阶导数
下界为$-2$、绝对值上界为$6$，又说明它同时是$2$弱凸且$6$光滑。这个例子不能推出
所有弱凸函数都满足PL条件，它只表明非凸性本身不排斥目标值的梯度控制。

式~\eqref{eq:opt-regularized-least-squares}的Hessian矩阵为
$\bm{X}^{\top}\bm{X}/n+\lambda\bm{I}$。当$\lambda>0$时，正则项把所有缺失曲率的
方向至少抬高到$\lambda$，因而目标$\lambda$强凸。若$\lambda=0$且$\bm{X}$不满列秩，
驻点仍是全局最优点，但解不唯一；若目标再加入非凸项，驻点甚至不再保证局部最小。

\section{梯度下降与步长}
\label{sec:opt-gradient-descent}

\subsection{光滑上界与单步下降}

梯度给出最陡的一阶下降方向，但一步走多远取决于函数弯曲程度。光滑性把局部线性近似的
余项控制为二次量\scite{math-nocedal2006numerical}

\begin{lemma}[下降引理]
\label{lem:opt-descent}
若$f$的梯度是$L$-Lipschitz连续的，则
\begin{equation}
  f(\bm{v})\leq f(\bm{u})
  +\langle\nabla f(\bm{u}),\bm{v}-\bm{u}\rangle
  +\frac{L}{2}\lVert\bm{v}-\bm{u}\rVert_2^2.
  \label{eq:opt-descent-lemma}
\end{equation}
\end{lemma}

\begin{proof}
令$\bm{d}=\bm{v}-\bm{u}$。沿线段积分并减去一阶项可得
\[
  f(\bm{v})-f(\bm{u})-\langle\nabla f(\bm{u}),\bm{d}\rangle
  =
  \int_0^1
  \langle\nabla f(\bm{u}+t\bm{d})-\nabla f(\bm{u}),\bm{d}\rangle\,dt.
\]
由Cauchy--Schwarz不等式和梯度Lipschitz条件，被积函数至多为
$Lt\lVert\bm{d}\rVert_2^2$；积分即得结论。
\end{proof}

梯度下降按
\begin{equation}
  \bm{w}^{(t+1)}
  =\bm{w}^{(t)}-\eta\nabla f(\bm{w}^{(t)})
  \label{eq:opt-gradient-descent}
\end{equation}
更新。把$\bm{v}$取为右端点代入下降引理，得到
\begin{equation}
  f(\bm{w}^{(t+1)})
  \leq f(\bm{w}^{(t)})
  -\eta\left(1-\frac{L\eta}{2}\right)
  \lVert\nabla f(\bm{w}^{(t)})\rVert_2^2.
  \label{eq:opt-gradient-single-step}
\end{equation}
所以$0<\eta<2/L$保证每个非驻点处严格下降。若$f$下有界，记其有限下确界为
\[
  f_{\inf}:=\inf_{\bm{w}\in\mathbb{R}^d}f(\bm{w})>-\infty.
\]
取$\eta=1/L$并对$t=0,\ldots,T-1$求和，可得非凸光滑情形的保证
\begin{equation}
  \min_{0\leq t<T}\lVert\nabla f(\bm{w}^{(t)})\rVert_2^2
  \leq
  \frac{2L(f(\bm{w}^{(0)})-f_{\inf})}{T}.
  \label{eq:opt-nonconvex-gradient-rate}
\end{equation}
它只说明某次迭代接近驻点，没有声称驻点是全局最优点。

步长上界在一维二次函数上即可看到。对$f(w)=Lw^2/2$，
$w^{(t+1)}=(1-\eta L)w^{(t)}$。当$\eta=2/L$时迭代在正负之间等幅振荡；
$\eta>2/L$时幅度按$|1-\eta L|>1$增长。负梯度方向本身没有错，失败来自步长跨过了
曲率允许的稳定范围。

\subsection{不同曲率下的收敛证明}

若$f$还凸、最小点$\bm{w}^{\star}$存在，并继续取$\eta=1/L$，则由凸性有
$f(\bm{w}^{(t)})-f^\star
\leq\langle\nabla f(\bm{w}^{(t)}),\bm{w}^{(t)}-\bm{w}^{\star}\rangle$。
先展开距离平方并使用这个不等式：
\begin{align*}
  \lVert\bm{w}^{(t+1)}-\bm{w}^{\star}\rVert_2^2
  &=
  \lVert\bm{w}^{(t)}-\bm{w}^{\star}\rVert_2^2
  -\frac{2}{L}
  \left\langle\nabla f(\bm{w}^{(t)}),
  \bm{w}^{(t)}-\bm{w}^{\star}\right\rangle\\
  &\qquad
  +\frac{1}{L^2}\lVert\nabla f(\bm{w}^{(t)})\rVert_2^2\\
  &\leq
  \lVert\bm{w}^{(t)}-\bm{w}^{\star}\rVert_2^2
  -\frac{2}{L}\bigl(f(\bm{w}^{(t)})-f^\star\bigr)\\
  &\qquad
  +\frac{1}{L^2}\lVert\nabla f(\bm{w}^{(t)})\rVert_2^2.
\end{align*}
另一方面，把$\eta=1/L$代入式~\eqref{eq:opt-gradient-single-step}可得
\[
  \frac{1}{L^2}\lVert\nabla f(\bm{w}^{(t)})\rVert_2^2
  \leq
  \frac{2}{L}\bigl(
  f(\bm{w}^{(t)})-f(\bm{w}^{(t+1)})\bigr).
\]
用后一式消去梯度平方项，得到
\[
  \lVert\bm{w}^{(t+1)}-\bm{w}^{\star}\rVert_2^2
  \leq
  \lVert\bm{w}^{(t)}-\bm{w}^{\star}\rVert_2^2
  -\frac{2}{L}\bigl(f(\bm{w}^{(t+1)})-f^\star\bigr).
\]
对$t=0,\ldots,T-1$求和并去掉末端的非负距离，便有
\[
  \frac{2}{L}\sum_{t=0}^{T-1}
  \bigl(f(\bm{w}^{(t+1)})-f^\star\bigr)
  \leq
  \lVert\bm{w}^{(0)}-\bm{w}^{\star}\rVert_2^2.
\]
式~\eqref{eq:opt-gradient-single-step}还说明目标值单调不增，因此左侧求和中的每一项
都不小于$f(\bm{w}^{(T)})-f^\star$。于是
\begin{equation}
  f(\bm{w}^{(T)})-f^\star
  \leq\frac{L\lVert\bm{w}^{(0)}-\bm{w}^{\star}\rVert_2^2}{2T}.
  \label{eq:opt-convex-gradient-rate}
\end{equation}

若满足PL条件，把式~\eqref{eq:opt-pl-condition}代入
式~\eqref{eq:opt-gradient-single-step}，在$\eta=1/L$时得到
\begin{equation}
  f(\bm{w}^{(t)})-f^\star
  \leq
  \left(1-\frac{\mu}{L}\right)^t
  \bigl(f(\bm{w}^{(0)})-f^\star\bigr).
  \label{eq:opt-pl-linear-rate}
\end{equation}
强凸函数满足PL条件，还由
$f(\bm{w})-f^\star\geq\mu\lVert\bm{w}-\bm{w}^{\star}\rVert_2^2/2$
把目标误差转成参数误差。没有强凸时，式~\eqref{eq:opt-pl-linear-rate}不保证迭代趋向
某个唯一参数。

固定步长需要已知或估计$L$。回溯线搜索从候选步长开始缩小，直到
$f(\bm{w}-\eta\nabla f(\bm{w}))\leq
f(\bm{w})-c\eta\lVert\nabla f(\bm{w})\rVert_2^2$成立，其中$c\in(0,1)$。
学习率调度则预先或按进度改变$\eta_t$。线搜索用实际函数值核验当前一步，调度只规定时间
表；二者都不能弥补错误梯度或无下界目标。

\section{曲率、加速与预条件}
\label{sec:opt-curvature-acceleration}

\subsection{动量与加速的递推结构}

对二次函数$f(\bm{w})=\frac12\bm{w}^{\top}\bm{H}\bm{w}
-\bm{c}^{\top}\bm{w}$，梯度下降的误差满足
$\bm{e}^{(t+1)}=(\bm{I}-\eta\bm{H})\bm{e}^{(t)}$。当
$\bm{H}\succ\bm{0}$但条件数$\kappa=L/\mu$很大时，同一步长必须同时照顾陡峭与平坦
方向，平坦方向因而收敛缓慢。

\term{动量法}（momentum method）保存过去方向：
\[
  \bm{v}^{(t+1)}=\beta\bm{v}^{(t)}+\nabla f(\bm{w}^{(t)}),
  \qquad
  \bm{w}^{(t+1)}=\bm{w}^{(t)}-\eta\bm{v}^{(t+1)}.
\]
它通过二阶递推重塑各特征方向上的多项式响应，但$\beta$和$\eta$不能各自独立地任取。
对上述正定二次目标，曲率为$\lambda_i$的特征方向满足
\[
  e_i^{(t+1)}
  =(1+\beta-\eta\lambda_i)e_i^{(t)}-\beta e_i^{(t-1)}.
\]
该方向的两个特征根都严格位于单位圆内，当且仅当
$-1<\beta<1$且$0<\eta\lambda_i<2(1+\beta)$。因而常用的
$0\leq\beta<1$还须配合$0<\eta<2(1+\beta)/L$，才能保证所有方向稳定。
若已知$0<\mu\leq\lambda_i\leq L$，经典重球参数
\[
  \eta=\frac{4}{(\sqrt{L}+\sqrt{\mu})^2},
  \qquad
  \beta=\left(\frac{\sqrt{L}-\sqrt{\mu}}
  {\sqrt{L}+\sqrt{\mu}}\right)^2
\]
给出这一二次模型上的代表性最坏方向速率；它不是一般非凸目标或含噪梯度下的无条件
调参规则。

\term{Nesterov加速梯度}
（Nesterov accelerated gradient）在前瞻点计算梯度：
\[
  \bm{z}^{(t)}=\bm{w}^{(t)}+\beta_t(\bm{w}^{(t)}-\bm{w}^{(t-1)}),
  \qquad
  \bm{w}^{(t+1)}=\bm{z}^{(t)}-\frac1L\nabla f(\bm{z}^{(t)}).
\]
令$\bm{w}^{(-1)}=\bm{w}^{(0)}$、$q_0=1$，并取
\[
  q_{t+1}=\frac{1+\sqrt{1+4q_t^2}}{2},
  \qquad
  \beta_t=\frac{q_t-1}{q_{t+1}}.
\]
若$f$为凸可微函数、梯度$L$-Lipschitz连续、最小点存在，并且每步使用精确梯度和
步长$1/L$，则上述选择给出
$f(\bm{w}^{(t)})-f^\star=O(1/t^2)$；同样条件下，普通梯度下降的一般保证为
$O(1/t)$。强凸时还可采用依赖$\mu/L$的固定系数得到线性速率，但那是另一组参数和证明。
在一维二次函数$f(w)=\lambda w^2/2$上，这个迭代可直接核验为
\[
  w^{(t+1)}
  =\left(1-\frac{\lambda}{L}\right)
  \bigl((1+\beta_t)w^{(t)}-\beta_t w^{(t-1)}\bigr).
\]
当$\lambda=L$时一步到达零点；当$\lambda\ll L$时，跨步项改变平坦方向上的二阶递推，
但系数选错仍会振荡。一般凸证明不能只检查这个二次算例，而是构造同时含目标误差与外推
距离的势函数；$q_{t+1}^2-q_{t+1}=q_t^2$使相邻两步的耦合项抵消，最终留下
$q_t^2(f(\bm{w}^{(t)})-f^\star)$的统一上界。加速结论因此来自特定外推序列与势函数
不等式，不能由“加入惯性”本身推出。

\subsection{牛顿、近似曲率与正定预条件}

\term{牛顿法}（Newton's method）在$\bm{w}^{(t)}$处使用关于增量$\bm{p}$的局部
二次模型
\[
  m_t(\bm{p})
  =
  f(\bm{w}^{(t)})
  +\left\langle\nabla f(\bm{w}^{(t)}),\bm{p}\right\rangle
  +\frac12\bm{p}^{\top}\nabla^2f(\bm{w}^{(t)})\bm{p}.
\]
令$\nabla_{\bm{p}}m_t(\bm{p}^{(t)})=\bm{0}$，便得到牛顿方程
\[
  \nabla^2 f(\bm{w}^{(t)})\bm{p}^{(t)}
  =-\nabla f(\bm{w}^{(t)}),\qquad
  \bm{w}^{(t+1)}=\bm{w}^{(t)}+\alpha_t\bm{p}^{(t)}.
\]
若Hessian矩阵在最优点附近正定且Lipschitz连续，取单位步长可得到局部二次收敛。但Hessian矩阵
不定时，牛顿方向未必下降；离最优点较远时还需线搜索或信赖域。

对非线性最小二乘
$f(\bm{w})=\frac12\lVert\bm{r}(\bm{w})\rVert_2^2$，设
$\bm{r}:\mathbb{R}^d\to\mathbb{R}^m$。记残差向量的Jacobian矩阵为
$\bm{J}(\bm{w})\in\mathbb{R}^{m\times d}$，其第$i$行为
$\bm{J}_{i,:}(\bm{w})=\nabla r_i(\bm{w})^{\top}$，则
\[
  \nabla^2f(\bm{w})
  =\bm{J}(\bm{w})^{\top}\bm{J}(\bm{w})
  +\sum_{i=1}^m r_i(\bm{w})\nabla^2r_i(\bm{w}).
\]
\term{Gauss--Newton法}忽略第二项，使用半正定矩阵
$\bm{J}^{\top}\bm{J}$。当残差小或残差函数近似线性时近似合理；残差大时，省略项可能
决定曲率。

\term{拟牛顿法}（quasi-Newton method）不显式计算Hessian矩阵，而用梯度差
$\bm{y}^{(t)}=\nabla f(\bm{w}^{(t+1)})-\nabla f(\bm{w}^{(t)})$和步长
$\bm{s}^{(t)}=\bm{w}^{(t+1)}-\bm{w}^{(t)}$更新矩阵近似，并满足割线条件
$\bm{B}^{(t+1)}\bm{s}^{(t)}=\bm{y}^{(t)}$。若$\bm{B}^{(t)}$近似Hessian矩阵，
BFGS更新为
\begin{equation}
  \bm{B}^{(t+1)}
  =
  \bm{B}^{(t)}
  +
  \frac{\bm{y}^{(t)}(\bm{y}^{(t)})^{\top}}
       {(\bm{y}^{(t)})^{\top}\bm{s}^{(t)}}
  -
  \frac{\bm{B}^{(t)}\bm{s}^{(t)}(\bm{s}^{(t)})^{\top}\bm{B}^{(t)}}
       {(\bm{s}^{(t)})^{\top}\bm{B}^{(t)}\bm{s}^{(t)}}.
  \label{eq:opt-bfgs-update}
\end{equation}
当$\bm{B}^{(t)}\succ\bm{0}$且曲率条件
$(\bm{y}^{(t)})^{\top}\bm{s}^{(t)}>0$成立时，
式~\eqref{eq:opt-bfgs-update}保持$\bm{B}^{(t+1)}\succ\bm{0}$，相应搜索方向
$-(\bm{B}^{(t+1)})^{-1}\nabla f$才由正定性保证是下降方向。沿下降方向执行满足
Wolfe条件的线搜索可产生该正曲率内积。具体地，令
$\bm{w}^{(t+1)}=\bm{w}^{(t)}+\alpha_t\bm{p}^{(t)}$，其中
$\nabla f(\bm{w}^{(t)})^{\top}\bm{p}^{(t)}<0$。Wolfe条件要求存在
$0<c_1<c_2<1$使
\begin{equation}
\begin{aligned}
  f(\bm{w}^{(t+1)})
  &\leq
  f(\bm{w}^{(t)})
  +
  c_1\alpha_t
  \nabla f(\bm{w}^{(t)})^{\top}\bm{p}^{(t)},\\
  \nabla f(\bm{w}^{(t+1)})^{\top}\bm{p}^{(t)}
  &\geq
  c_2\nabla f(\bm{w}^{(t)})^{\top}\bm{p}^{(t)}.
\end{aligned}
\label{eq:opt-wolfe-conditions}
\end{equation}
第二行就是\term{Wolfe曲率条件}（Wolfe curvature condition）。由
$\bm{s}^{(t)}=\alpha_t\bm{p}^{(t)}$可得
\[
  (\bm{y}^{(t)})^{\top}\bm{s}^{(t)}
  \geq
  \alpha_t(c_2-1)
  \nabla f(\bm{w}^{(t)})^{\top}\bm{p}^{(t)}
  >0.
\]
非凸或梯度含噪时，上述条件可能难以满足，此时可能需要阻尼或跳过更新。

\term{预条件}（preconditioning）则选择正定矩阵$\bm{P}$，按
\[
  \bm{w}^{(t+1)}
  =
  \bm{w}^{(t)}-\eta\bm{P}\nabla f(\bm{w}^{(t)})
\]
更新。
这等价于在新的坐标或内积下做最陡下降。若\(p_i>0\)，并且
\[
  \bm{P}
  =
  \bm{Q}\operatorname{diag}(p_1,\ldots,p_d)\bm{Q}^{\top},
\]
则定义
\[
  \bm{P}^{1/2}
  =
  \bm{Q}\operatorname{diag}(\sqrt{p_1},\ldots,\sqrt{p_d})\bm{Q}^{\top}.
\]
这就是$\bm{P}$唯一的对称正定平方根：它满足
$\bm{P}^{1/2}\bm{P}^{1/2}=\bm{P}$，并且可逆，其逆记为$\bm{P}^{-1/2}$。
对当前正定二次目标，记
\[
  \widetilde{\bm{H}}=\bm{P}^{1/2}\bm{H}\bm{P}^{1/2},
  \qquad
  0<\widetilde{\mu}
  =\lambda_{\min}(\widetilde{\bm{H}})
  \leq\lambda_{\max}(\widetilde{\bm{H}})
  =\widetilde{L}.
\]
原坐标中的误差递推为
$\bm{e}^{(t+1)}=(\bm{I}-\eta\bm{P}\bm{H})\bm{e}^{(t)}$。令
$\widetilde{\bm{e}}^{(t)}=\bm{P}^{-1/2}\bm{e}^{(t)}$，则
\[
  \widetilde{\bm{e}}^{(t+1)}
  =
  (\bm{I}-\eta\widetilde{\bm{H}})
  \widetilde{\bm{e}}^{(t)},
  \qquad
  \bm{P}^{-1/2}
  (\bm{I}-\eta\bm{P}\bm{H})
  \bm{P}^{1/2}
  =
  \bm{I}-\eta\widetilde{\bm{H}}.
\]
右侧矩阵由可逆坐标变换$\bm{P}^{1/2}$得到，因而与原误差递推矩阵相似；相似矩阵具有
相同特征值。由于$\widetilde{\bm{H}}$对称正定，迭代对任意初始误差收敛当且仅当
$0<\eta<2/\widetilde{L}$。取
$\eta=2/(\widetilde{L}+\widetilde{\mu})$时，在变换后的欧氏范数中，每步最坏收缩因子为
\[
  \frac{\widetilde{\kappa}-1}{\widetilde{\kappa}+1},
  \qquad
  \widetilde{\kappa}
  =\frac{\widetilde{L}}{\widetilde{\mu}}.
\]
因此预条件的目标不是单独让$\bm{P}$的条件数小，而是让变换后曲率
$\widetilde{\bm{H}}$的条件数接近$1$。固定$\bm{P}\succ\bm{0}$只保证定义了合法内积；
若它与$\bm{H}$不匹配，$\widetilde{\kappa}$反而可能增大。对一般二阶可微的光滑强凸
函数，需要在所有迭代经过的区域统一满足
$\widetilde{\mu}\bm{I}\preceq
\bm{P}^{1/2}\nabla^2f(\bm{w})\bm{P}^{1/2}
\preceq\widetilde{L}\bm{I}$，上述步长范围与线性收敛解释才相应延续；变化的预条件矩阵还需
额外控制。

动量利用历史更新，牛顿类近似当前曲率，预条件改变几何尺度；三者可能组合，却不是同一种
机制。第~\ref{chap:matrix-analysis}章将在矩阵函数与Kronecker结构建立后给出结构化
预条件，Fisher矩阵与经验梯度二阶矩的差别则需等概率对象建立后讨论。

\section{约束与拉格朗日乘子}
\label{sec:opt-kkt}

\subsection{可行方向与KKT条件}

资源约束使最优点处的目标梯度不必为零，因为可行方向受到限制。对
\[
  \min_{\bm{w}}f(\bm{w})
  \quad\text{s.t.}\quad
  a_j(\bm{w})=0\ (j=1,\ldots,m),\quad
  g_k(\bm{w})\leq0\ (k=1,\ldots,p),
\]
在可行点$\bm{w}$处，满足
\[
  \langle\nabla a_j(\bm{w}),\bm{d}\rangle=0,\qquad
  \langle\nabla g_k(\bm{w}),\bm{d}\rangle\leq0
  \quad\text{对所有活跃的 }g_k
\]
的方向$\bm{d}$称为\term{一阶可行方向}（linearized feasible direction）。
等式约束为何产生法向量平衡，可以先看仿射情形
$\bm{A}\bm{w}=\bm{b}$。任一可行扰动都满足$\bm{A}\bm{d}=\bm{0}$；局部最优点
不能沿$\bm{d}$或$-\bm{d}$下降，故
$\langle\nabla f(\bm{w}^{\star}),\bm{d}\rangle=0$对所有
$\bm{d}\in\ker\bm{A}$成立。由
$(\ker\bm{A})^\perp=\operatorname{range}(\bm{A}^{\top})$，存在不受符号限制的
$\bm{\nu}$使
$\nabla f(\bm{w}^{\star})=-\bm{A}^{\top}\bm{\nu}$。

活跃不等式只允许切平面一侧的一阶移动，它的外法向量因而以非负系数参与梯度平衡；
不活跃约束在该点留有严格余量，不限制足够小的移动，对应系数应为零。这分别引出
非负乘子与互补松弛。在适当的约束资格条件下，这些几何关系可通过
\term{拉格朗日函数}（Lagrangian）统一表达：
\[
  \mathcal{L}(\bm{w},\bm{\nu},\bm{\alpha})
  =f(\bm{w})+\sum_j\nu_j a_j(\bm{w})
  +\sum_k\alpha_k g_k(\bm{w}),
  \qquad
  \bm{\nu}\in\mathbb{R}^m,\quad
  \bm{\alpha}\in\mathbb{R}_+^p.
\]
乘子把违反约束的边际代价并入目标。

\term{Karush--Kuhn--Tucker条件}（KKT conditions）包括
\begin{align}
  \nabla_{\bm{w}}\mathcal{L}(\bm{w}^{\star},\bm{\nu}^{\star},
  \bm{\alpha}^{\star})&=\bm{0},
  &&\text{驻点性},\label{eq:opt-kkt-stationarity}\\
  a_j(\bm{w}^{\star})=0,\quad g_k(\bm{w}^{\star})&\leq0,
  &&\text{原始可行性},\\
  \alpha_k^{\star}&\geq0,
  &&\text{对偶可行性},\\
  \alpha_k^{\star}g_k(\bm{w}^{\star})&=0,
  &&\text{互补松弛}.\label{eq:opt-kkt-complementarity}
\end{align}
互补松弛表示未用尽的约束不能有正价格，正价格只可能出现在活跃边界上。

局部最优点满足KKT条件需要\term{约束资格条件}（constraint qualification）。
例如活跃不等式梯度与等式梯度线性无关时，沿可行方向的一阶下降不可能存在，由有限维分离
可得到乘子。没有资格条件，乘子可能不存在。例$f(x)=x$且约束$x^2\leq0$只有可行点
$x=0$，它当然最优；但约束梯度为零，驻点条件$1+\alpha\cdot0=0$无解。

反方向在凸情形更强。若$f$和每个$g_k$均为可微凸函数，$a_j$仿射，且某个可行点与乘子
满足KKT，则对任意可行$\bm{w}$，仿射等式约束使
$\langle\nabla a_j(\bm{w}^{\star}),\bm{w}-\bm{w}^{\star}\rangle
=a_j(\bm{w})-a_j(\bm{w}^{\star})=0$。因此等式乘子项消失，再由凸性和驻点性，
\[
  f(\bm{w})\geq
  f(\bm{w}^{\star})
  -\sum_k\alpha_k^{\star}
  \langle\nabla g_k(\bm{w}^{\star}),\bm{w}-\bm{w}^{\star}\rangle
  \geq f(\bm{w}^{\star}).
\]
最后一步使用凸约束的一阶下界、可行性和互补松弛。因此KKT在该条件下足以保证全局最优。

\subsection{活跃集与解的敏感性}

\begin{example}[预算约束、活跃集与敏感性]
\label{ex:opt-kkt-budget}
求最接近$\bm{c}=(2,1)^{\top}$的非负分配，并令总预算为$\tau>0$：
\[
  \min_{\bm{w}}\frac12\lVert\bm{w}-\bm{c}\rVert_2^2
  \quad\text{s.t.}\quad
  w_1+w_2\leq\tau,\quad -w_1\leq0,\quad -w_2\leq0.
\]
当$\tau>3$时，$\bm{c}$本身严格满足预算约束，故该约束不活跃，且
$\bm{w}^{\star}(\tau)=(2,1)^{\top}$。当$\tau=3$时解仍为$\bm{c}$，预算约束活跃，
但其乘子为零；这个边界说明“约束活跃”不等于“乘子严格为正”。
当$1<\tau<3$时，两个坐标仍为正而预算取等。把预算这一不等式的乘子记为
$\alpha_{\mathrm{bud}}$，驻点条件给
\[
  w_1-2+\alpha_{\mathrm{bud}}=0,\qquad
  w_2-1+\alpha_{\mathrm{bud}}=0,\qquad
  w_1+w_2=\tau.
\]
因此
\[
  \bm{w}^{\star}(\tau)
  =\left(\frac{\tau+1}{2},\frac{\tau-1}{2}\right)^{\top},
  \qquad
  \alpha_{\mathrm{bud}}^{\star}(\tau)=\frac{3-\tau}{2}.
\]
当$0<\tau<1$时，第二个非负约束也活跃，解变成
$\bm{w}^{\star}(\tau)=(\tau,0)^{\top}$，预算乘子为$2-\tau$，第二个非负约束的
乘子为$1-\tau$。在$\tau=1$处，第二个非负约束活跃但乘子为零。上述区域及边界分别
展示预算不生效、约束活跃但零乘子、仅预算产生正价格，以及预算与坐标边界同时产生
正价格。目标严格凸，满足KKT的点都是唯一全局最优解。
\end{example}

这个算例也把两类敏感性分开。在$1<\tau<3$时，
$d\bm{w}^{\star}/d\tau=(1/2,1/2)^{\top}$；在$0<\tau<1$时则为
$(1,0)^{\top}$。$\tau=1$处活跃集切换，解连续但导数改变。最优值为
\[
  V(\tau)=
  \begin{cases}
    0, & \tau\geq3,\\
    (3-\tau)^2/4, & 1\leq\tau\leq3,\\
    \bigl((2-\tau)^2+1\bigr)/2, & 0<\tau\leq1.
  \end{cases}
\]
在每个活跃集固定的开区间内，
$V'(\tau)=-\alpha_{\mathrm{bud}}^{\star}(\tau)$：放宽一单位预算带来的
最优值下降率正是预算乘子的相反数。解导数是向量，乘子给出的是最优值导数，二者不能混用。

一般地，设问题依赖参数$\tau$，活跃集在邻域内不变，把驻点条件、等式和活跃约束写成
$\bm{\Phi}(\bm{w},\bm{\nu},\bm{\alpha},\tau)=\bm{0}$。若对前三组变量的Jacobian
矩阵可逆，则第~\ref{chap:mathematical-analysis}章的隐式求导给出
\[
  \frac{d}{d\tau}
  \begin{bmatrix}\bm{w}^{\star}\\\bm{\nu}^{\star}\\\bm{\alpha}^{\star}\end{bmatrix}
  =
  -\left(\frac{\partial\bm{\Phi}}
  {\partial(\bm{w},\bm{\nu},\bm{\alpha})}\right)^{-1}
  \frac{\partial\bm{\Phi}}{\partial\tau}.
\]
这个Jacobian是完整KKT系统的导数，原目标Hessian可逆并不足以保证它可逆。活跃集改变或
该矩阵奇异时，解映射可能不可微；此时不能跨越切换点沿用同一个导数。

\section{原问题与对偶问题}
\label{sec:opt-duality}

\subsection{Lagrangian下界与弱对偶}

对固定乘子$\bm{\alpha}\geq\bm{0}$，任意可行$\bm{w}$都满足
$\mathcal{L}(\bm{w},\bm{\nu},\bm{\alpha})\leq f(\bm{w})$。因此
\term{对偶函数}（dual function）
\[
  q(\bm{\nu},\bm{\alpha})
  =\inf_{\bm{w}}\mathcal{L}(\bm{w},\bm{\nu},\bm{\alpha})
\]
给出原问题最优值的下界。最大化这个下界得到对偶问题。由构造立即有
$d^{\star}\leq p^{\star}$，称为\term{弱对偶}（weak duality）；差
$p^{\star}-d^{\star}$称为对偶间隙。对某组乘子取下确界得到$-\infty$，表示这组
乘子没有给出有限下界；对偶不可行则表示没有任何乘子满足对偶问题规定的可行条件。前者
描述一个候选乘子的函数值，后者描述整个候选集合，二者不能混用。

\subsection{Farkas择一与线性规划强对偶}

线性规划的强对偶可以从一个不可行证书推出。先建立所需的择一定理
\scite{math-boyd2004convex}

\begin{lemma}[有限生成锥的闭性]
\label{lem:opt-finitely-generated-cone}
给定$\bm{a}_1,\ldots,\bm{a}_r\in\mathbb{R}^m$，集合
\[
  K=\left\{\sum_{i=1}^r\alpha_i\bm{a}_i:\alpha_i\geq0\right\}
\]
是闭凸锥。
\end{lemma}

\begin{proof}
凸性和锥性质由非负线性组合直接得到。为证明闭性，先注意每个$\bm{x}\in K$都能由一组
线性无关的生成向量作非负组合。若某个表示使用的生成向量线性相关，就存在非零系数
$\bm{\gamma}$使$\sum_i\gamma_i\bm{a}_i=\bm{0}$；必要时改变符号，使至少一个
$\gamma_i>0$。从原系数中减去
$t\bm{\gamma}$，其中$t=\min_{\gamma_i>0}\alpha_i/\gamma_i$，仍保持所有系数
非负，并使至少一个系数变为零。重复此过程便得到线性无关的表示。

现取$\bm{x}^{(k)}\in K$且$\bm{x}^{(k)}\to\bm{x}$。线性无关的生成子集只有有限多个，
故可取子列，使每个$\bm{x}^{(k)}$都使用同一个子集$J$。以这些向量为列组成满列秩矩阵
$\bm{A}_J$，相应非负系数满足
\[
  \bm{\alpha}^{(k)}
  =(\bm{A}_J^{\top}\bm{A}_J)^{-1}\bm{A}_J^{\top}\bm{x}^{(k)}
  \longrightarrow
  (\bm{A}_J^{\top}\bm{A}_J)^{-1}\bm{A}_J^{\top}\bm{x}.
\]
极限系数仍逐项非负，且$\bm{x}$等于$\bm{A}_J$乘该极限系数，故$\bm{x}\in K$。
\end{proof}

\begin{lemma}[有限维多面体的线性像]
\label{lem:opt-polyhedral-image}
有限维多面体在线性映射下的像仍是多面体，因而是闭集。
\end{lemma}

\begin{proof}
设多面体$P=\{\bm{u}:\bm{B}\bm{u}\leq\bm{d}\}$，线性映射为
$\bm{v}=\bm{T}\bm{u}$。其图像
\[
  \{(\bm{u},\bm{v}):\bm{B}\bm{u}\leq\bm{d},\
  \bm{v}=\bm{T}\bm{u}\}
\]
是多面体，而$\bm{T}P$是该图像消去$\bm{u}$后的坐标投影。对一个待消变量，
Fourier--Motzkin消元分别收集其系数为正、为负和为零的不等式；每一对正、负系数
不等式消去该变量，零系数不等式原样保留。所得有限组非严格线性不等式恰好描述投影。
逐个消去$\bm{u}$的坐标，便得到$\bm{T}P$的有限线性不等式表示。多面体是有限个闭
半空间的交，因此闭。
\end{proof}

\begin{theorem}[Farkas择一定理]
\label{thm:opt-farkas}
给定$\bm{M}\in\mathbb{R}^{m\times d}$和$\bm{r}\in\mathbb{R}^m$，下列两个系统
恰有一个可行：
\[
  \bm{M}\bm{x}=\bm{r},\quad\bm{x}\geq\bm{0};
  \qquad
  \bm{M}^{\top}\bm{z}\geq\bm{0},\quad
  \bm{r}^{\top}\bm{z}<0.
\]
\end{theorem}

\begin{proof}
两者不能同时成立，因为若都成立，则
$0>\bm{r}^{\top}\bm{z}
=\bm{x}^{\top}\bm{M}^{\top}\bm{z}\geq0$，矛盾。
若第一个系统不可行，$\bm{r}$不属于锥
$K=\{\bm{M}\bm{x}:\bm{x}\geq\bm{0}\}$。它由$\bm{M}$的有限个列向量生成，故由
引理~\ref{lem:opt-finitely-generated-cone}可知$K$是闭凸锥。对$K$与$\bm{r}$应用
定理~\ref{thm:opt-projection-separation}，得到$\bm{a}$和$c$，使
$\langle\bm{a},\bm{k}\rangle\leq c<\langle\bm{a},\bm{r}\rangle$。
由于$\bm{0}\in K$，有$c\geq0$；又因为$t\bm{k}\in K$对所有$t\geq0$成立，
必有$\langle\bm{a},\bm{k}\rangle\leq0$，否则取充分大的$t$便会超过$c$。
令$\bm{z}=-\bm{a}$，可得$\langle\bm{z},\bm{r}\rangle<0$且
$\langle\bm{z},\bm{k}\rangle\geq0$对所有$\bm{k}\in K$成立。特别地，
对$\bm{M}$的每个列向量成立，故
$\bm{M}^{\top}\bm{z}\geq\bm{0}$。第二个系统可行。
\end{proof}

\begin{theorem}[有限维线性规划强对偶]
\label{thm:opt-lp-strong-duality}
考虑标准型原问题与对偶问题
\[
  \min_{\bm{x}\geq\bm{0}}\bm{c}^{\top}\bm{x}
  \quad\text{s.t.}\quad\bm{A}\bm{x}=\bm{b},
  \qquad
  \max_{\bm{y}}\bm{b}^{\top}\bm{y}
  \quad\text{s.t.}\quad\bm{A}^{\top}\bm{y}\leq\bm{c}.
\]
若原问题可行且有限下确界为$p^\star$，则原问题与对偶问题都有最优解，且两个最优值
都等于$p^\star$。
\end{theorem}

\begin{proof}
原始可行域
$P=\{\bm{x}:\bm{A}\bm{x}=\bm{b},\ \bm{x}\geq\bm{0}\}$是非空多面体。由
引理~\ref{lem:opt-polyhedral-image}，其目标值集合
\[
  W=\{\bm{c}^{\top}\bm{x}:\bm{x}\in P\}
\]
是$\mathbb{R}$中的非空多面体。集合$W$下有界，且一维多面体包含自己的有限下端点，
故$p^\star=\inf W$属于$W$。因此存在原始可行点$\bm{x}^{\star}$满足
$\bm{c}^{\top}\bm{x}^{\star}=p^\star$，原问题达到其最优值。

对任意$\varepsilon>0$，系统
\[
  \bm{A}\bm{x}=\bm{b},\quad
  \bm{c}^{\top}\bm{x}+s=p^\star-\varepsilon,\quad
  \bm{x}\geq\bm{0},\ s\geq0
\]
不可行，否则会有目标值小于$p^\star$的原始可行点。对增广矩阵应用
定理~\ref{thm:opt-farkas}，存在$(\bm{y}_{\varepsilon},\gamma_\varepsilon)$满足
\[
  \bm{A}^{\top}\bm{y}_{\varepsilon}
  +\gamma_\varepsilon\bm{c}\geq\bm{0},\quad
  \gamma_\varepsilon\geq0,\quad
  \bm{b}^{\top}\bm{y}_{\varepsilon}
  +\gamma_\varepsilon(p^\star-\varepsilon)<0.
\]
取任一原始可行$\bar{\bm{x}}$并与第一式作内积可知$\gamma_\varepsilon>0$，否则会与
最后一式矛盾。令
$\widehat{\bm{y}}_\varepsilon=-\bm{y}_\varepsilon/\gamma_\varepsilon$，便有
$\bm{A}^{\top}\widehat{\bm{y}}_\varepsilon\leq\bm{c}$且
$\bm{b}^{\top}\widehat{\bm{y}}_\varepsilon>p^\star-\varepsilon$。

记对偶可行多面体为
$D=\{\bm{y}:\bm{A}^{\top}\bm{y}\leq\bm{c}\}$。上面的构造说明$D$非空，且其
目标值集合
\[
  V=\{\bm{b}^{\top}\bm{y}:\bm{y}\in D\}
\]
对每个$\varepsilon>0$都包含一个大于$p^\star-\varepsilon$的值。弱对偶说明$V$中的
每个值都不超过$p^\star$，所以$\sup V=p^\star$。由
引理~\ref{lem:opt-polyhedral-image}，$V$是$\mathbb{R}$中的多面体；非空的一维
多面体只能是闭区间、闭射线、单点或整条实线。由于$V$有有限上确界，它必包含上端点
$p^\star$。因此存在$\bm{y}^{\star}\in D$使
$\bm{b}^{\top}\bm{y}^{\star}=p^\star$，对偶最优解达到与原问题相同的最优值。
\end{proof}

\begin{example}[二维资源分配的原始与对偶证书]
\label{ex:opt-resource-lp-duality}
考虑
\[
  \max_{x_1,x_2\geq0}3x_1+2x_2
  \quad\text{s.t.}\quad
  2x_1+x_2\leq4,\quad x_1+2x_2\leq4.
\]
为两类资源设置价格$u_1,u_2\geq0$。由于原问题是最大化，对可行点的资源余量按非负
价格计价会给出目标上界：
\begin{align*}
  \mathcal{L}(\bm{x},\bm{u})
  &=
  3x_1+2x_2
  +u_1(4-2x_1-x_2)+u_2(4-x_1-2x_2)\\
  &=
  4u_1+4u_2
  +(3-2u_1-u_2)x_1+(2-u_1-2u_2)x_2.
\end{align*}
要使$\sup_{\bm{x}\geq\bm{0}}\mathcal{L}(\bm{x},\bm{u})$有限，$x_1,x_2$的系数
必须非正，因而
$2u_1+u_2\geq3$且$u_1+2u_2\geq2$。此时上确界就是$4u_1+4u_2$，
最小化这个上界得到对偶
\[
  \min_{u_1,u_2\geq0}4u_1+4u_2
  \quad\text{s.t.}\quad
  2u_1+u_2\geq3,\quad u_1+2u_2\geq2.
\]
原始点$(x_1,x_2)=(4/3,4/3)$可行，值为$20/3$；对偶点
$(u_1,u_2)=(4/3,1/3)$可行，值也为$20/3$。弱对偶说明二者同时最优。两项资源约束
和两项产品约束均取等号，互补松弛也逐项成立。
\end{example}

\subsection{一般凸问题中的资格条件}

不可行、无界和未达到必须区分。Farkas向量证明约束不可行；若原问题沿可行射线无限改善，
弱对偶只说明不存在给出有限下界的对偶候选，也就是每个允许乘子的对偶函数值都为
$-\infty$，不能据此断言乘子候选集合为空。在线性规划标准型中，这种情形表现为对偶
不可行。若只有下确界而不达到，则不能展示原始最优点。一般凸问题在使用Slater条件前还需
处理低维定义域。扩展实值函数$f$的\term{有效定义域}
（effective domain）是
$\operatorname{dom}f=\{\bm{w}:f(\bm{w})<+\infty\}$；
集合$C$的\term{仿射包}（affine hull）是包含$C$的最小仿射集合；
\term{相对内部}（relative interior）则是在$\operatorname{aff}C$的相对拓扑中具有
邻域的点。例如平面中的线段没有二维内部，但去掉两个端点后就是它的相对内部。

本章采用如下Slater版本。设$f$和不等式函数$g_k$都是适当凸函数，即它们不恒为
$+\infty$、从不取$-\infty$且至少在一点取有限值；这些函数具有共同凸定义域$D$，
等式约束为仿射约束$\bm{A}\bm{w}=\bm{b}$。若存在
\[
  \bar{\bm{w}}\in\operatorname{ri}D,\qquad
  \bm{A}\bar{\bm{w}}=\bm{b},\qquad
  g_k(\bar{\bm{w}})<0\quad\text{对所有 }k,
\]
并且原问题最优值有限，则原、对偶最优值相等，且存在最优对偶乘子。证明把可达到的
“约束余量—目标值”集合与严格低于原最优值的点分离；相对内部中的严格可行点排除了
目标系数为零的异常分离向量，归一化后便得到所需乘子。这里列出的定义域、仿射等式和
有限最优值都是所采用版本的一部分。凸性本身并不够，也不能把前面的线性规划证明原样
套到任意凸问题。第~\ref{chap:decision-theory-and-dynamic-risk}章会在状态、策略和
价值定义后，把同一对偶工具用于价值函数与占用度量。

\section{松弛与最优性证书}
\label{sec:opt-relaxations-certificates}

\subsection{凸包与可计算松弛}

离散约束使可行集不再凸。例如
\[
  S=\{(0,0),(1,0),(0,1)\}
\]
可以表示“至多选择一个项目”。把$S$替换为三角形
\[
  \operatorname{conv}(S)
  =\{(x_1,x_2):x_1,x_2\geq0,\ x_1+x_2\leq1\}
\]
得到\term{凸松弛}（convex relaxation）：原来的每个候选仍可行，但增加了分数候选。
对最小化问题，松弛最优值是原最优值的下界；对最大化问题则是上界。

精确凸包不改变任意线性目标$\bm{c}^{\top}\bm{x}$的最优值：任意
$\bm{x}=\sum_i\alpha_i\bm{s}_i\in\operatorname{conv}(S)$都满足
$\bm{c}^{\top}\bm{x}=\sum_i\alpha_i\bm{c}^{\top}\bm{s}_i
\geq\min_{\bm{s}\in S}\bm{c}^{\top}\bm{s}$，再结合
$S\subseteq\operatorname{conv}(S)$即得最小值相等，最大值同理。但它的完整不等式描述
可能难以获得。只保留局部约束通常
得到更大的多面体，计算较容易而界更松。也可以把乘积$x_ix_j$提升为矩阵变量并要求某个
块矩阵半正定，得到半正定松弛；它利用二次一致性，代价是矩阵锥约束更昂贵。

\subsection{从界到证书}

取整还可能破坏约束。考虑
\[
  \max_{x_1,x_2\in\{0,1\}}x_1+x_2
  \quad\text{s.t.}\quad 2x_1+2x_2\leq3.
\]
整数最优值为$1$，而把变量放宽到$[0,1]$后，$(3/4,3/4)$可行且目标值为$3/2$。
逐坐标四舍五入得到$(1,1)$，却违反容量约束；逐坐标向下取整虽可行，却得到目标值$0$。
因此松弛值、恢复可行性的规则和恢复后的目标损失必须分别分析。

可计算界构成停止证书。对最小化问题，若可行解给出上界$U$，对偶或松弛给出下界$L$，
则
\[
  0\leq U-p^\star\leq U-L.
\]
当$U=L$时已证明全局最优。若把松弛解取整为可行解$\widehat{\bm{x}}$，还应单独分析
$f(\widehat{\bm{x}})-L$；“取整后看起来接近”不是误差界。互补松弛能够证明哪些约束在
最优点活跃，却不总能直接构造整数解。图模型中的边缘多面体和MAP松弛会在图与概率对象
建立后复用这里的凸包、上下界和间隙概念。

\section{非光滑目标与正则化}
\label{sec:opt-nonsmooth-proximal}

\subsection{次梯度与非光滑最优性}

绝对值损失和$\ell_1$稀疏惩罚在零点不可微，但凸函数仍可用支撑直线描述。
向量$\bm{g}$若满足
\[
  f(\bm{v})\geq f(\bm{u})+\langle\bm{g},\bm{v}-\bm{u}\rangle
  \quad\text{对所有 }\bm{v},
\]
就称为$f$在$\bm{u}$处的\term{次梯度}（subgradient）。全部次梯度组成
$\partial f(\bm{u})$。例如
\[
  \partial |x|=
  \begin{cases}
    \{1\},&x>0,\\
    [-1,1],&x=0,\\
    \{-1\},&x<0.
  \end{cases}
\]
凸函数在$\bm{w}^{\star}$处最优，当且仅当
$\bm{0}\in\partial f(\bm{w}^{\star})$。若零向量是次梯度，定义立即给出
$f(\bm{v})\geq f(\bm{w}^{\star})$；反过来，若$\bm{w}^{\star}$全局最优，则
$f(\bm{v})\geq f(\bm{w}^{\star})+\langle\bm{0},\bm{v}-\bm{w}^{\star}\rangle$，
所以零向量满足次梯度定义。这个等价性依赖凸次梯度采用全局下界定义。
若要把$f+r$的次微分直接写成$\partial f+\partial r$，或把约束写成指标函数后拆分
次微分，还需共同相对内部等相应资格条件；不能只凭两个函数都凸就无条件使用求和等式。

\subsection{投影、邻近与软阈值}

对闭凸集$C$，欧氏\term{投影}（projection）
$\Pi_C(\bm{z})=\argmin_{\bm{w}\in C}\lVert\bm{w}-\bm{z}\rVert_2^2$
把一步梯度更新拉回可行域：
\[
  \bm{w}^{(t+1)}
  =\Pi_C\bigl(\bm{w}^{(t)}-\eta_t\nabla f(\bm{w}^{(t)})\bigr).
\]
定理~\ref{thm:opt-projection-separation}的证明已经说明，在有限维空间中，非空闭性保证
最近点存在，而凸性与平方距离的严格凸性共同保证唯一。若$C$非凸，最近点可能不唯一。

更一般地，函数$r$的\term{邻近算子}（proximal operator）定义为
\[
  \operatorname{prox}_{\eta r}(\bm{z})
  =\argmin_{\bm{w}}
  \left\{r(\bm{w})+\frac{1}{2\eta}
  \lVert\bm{w}-\bm{z}\rVert_2^2\right\}.
\]
这里取$\eta>0$，并假设$r$是前述意义下的适当函数、下半连续且凸。二次项使目标强凸，
强制性与下半连续性保证最小点存在，强凸性保证它唯一，所以邻近算子对每个$\bm{z}$都有
单值输出。
指标函数$\iota_C$在$C$上为零、在外部为$+\infty$，其邻近算子正是投影。
对$f+r$，近端梯度先按光滑部分$f$走一步，再用$r$的邻近算子处理非光滑结构。

\begin{theorem}[$\ell_1$邻近算子是软阈值]
\label{thm:opt-soft-threshold}
对$r(\bm{w})=\lambda\lVert\bm{w}\rVert_1$，
\[
  [\operatorname{prox}_{\eta r}(\bm{z})]_j
  =\operatorname{sign}(z_j)\max\{|z_j|-\eta\lambda,0\}.
\]
\end{theorem}

\begin{proof}
目标按坐标可分，只需最小化
$\lambda|w|+(w-z)^2/(2\eta)$。若$w>0$，驻点为
$w=z-\eta\lambda$，只在$z>\eta\lambda$时落在该区域；若$w<0$，驻点为
$w=z+\eta\lambda$，只在$z<-\eta\lambda$时有效。若
$|z|\leq\eta\lambda$，零点最优性条件
$0\in\lambda[-1,1]-z/\eta$成立，所以$w=0$。合并三种情形即得结论。
\end{proof}

邻近算子不会放大输入距离。令
$\bm{p}=\operatorname{prox}_{\eta r}(\bm{x})$、
$\bm{q}=\operatorname{prox}_{\eta r}(\bm{y})$。最优性条件给出
$(\bm{x}-\bm{p})/\eta\in\partial r(\bm{p})$和
$(\bm{y}-\bm{q})/\eta\in\partial r(\bm{q})$。把两个次梯度不等式相加，得到
\[
  \langle\bm{x}-\bm{y},\bm{p}-\bm{q}\rangle
  \geq\lVert\bm{p}-\bm{q}\rVert_2^2.
\]
再用Cauchy--Schwarz不等式便有
$\lVert\bm{p}-\bm{q}\rVert_2\leq\lVert\bm{x}-\bm{y}\rVert_2$。
这个非扩张性质依赖凸次梯度的单调性；非凸邻近子问题可能多解，不能直接沿用。

本节的凸次梯度也不同于第~\ref{chap:mathematical-analysis}章介绍的Clarke广义梯度。
前者用一个全局仿射下界定义，直接产生凸最优性条件；后者取局部Lipschitz函数邻近可微点
梯度极限的闭凸包，允许非凸局部结构。人为指定的替代梯度更只是反向计算信号，未必属于
前面任一集合。三者名称相近，定义对象与可用结论并不相同。

\subsection{惩罚、约束与参数表达}

惩罚与硬约束有联系，但不是逐参数自动等价。考虑约束问题
\[
  \min_{\bm{w}}f(\bm{w})
  \quad\text{s.t.}\quad
  \lVert\bm{w}\rVert_1\leq R.
\]
若该问题有原始最优解，强对偶成立，并且对偶最优乘子
$\lambda^\star$确实达到，例如满足前述Slater条件且最优值有限，则原始最优解同时最小化
$f(\bm{w})+\lambda^\star\lVert\bm{w}\rVert_1$；Lagrangian中的常数
$-\lambda^\star R$不影响最小点。仅有原、对偶最优值相等并不能保证这样的有限乘子存在。
反过来，不同$\lambda$对应的范数半径可能相同，非凸问题还可能存在对偶间隙。把
$\bm{w}$写成其他参数的函数也会改变优化几何；即使表示同一预测函数，也不保证产生同一
更新轨迹。

同样的区别出现在权重衰减中。普通欧氏梯度下降用于
$f(\bm{w})+\lambda\lVert\bm{w}\rVert_2^2/2$时给出
\[
  \bm{w}^{(t+1)}
  =(1-\eta\lambda)\bm{w}^{(t)}-\eta\nabla f(\bm{w}^{(t)}),
\]
与直接乘上衰减因子相同。若使用预条件矩阵$\bm{P}$，惩罚梯度产生
$-\eta\lambda\bm{P}\bm{w}^{(t)}$，而坐标无关的直接衰减仍为
$-\eta\lambda\bm{w}^{(t)}$；除非$\bm{P}=\bm{I}$或另有匹配，它们不再相同。

\section{随机梯度与更新稳定性}
\label{sec:opt-stochastic-gradient-stability}

\subsection{有限抽样与带误差迭代}

有限和目标
\[
  F_S(\bm{w})=\frac1n\sum_{i=1}^n\ell_i(\bm{w})
  +\frac{\lambda}{2}\lVert\bm{w}\rVert_2^2
\]
的经验损失梯度要读取全部$n$项。若$I$在$\{1,\ldots,n\}$上均匀抽取，则完整目标的
随机梯度可取
\[
  \bm{g}_I(\bm{w})
  =\nabla\ell_I(\bm{w})+\lambda\bm{w},\qquad
  \mathbb{E}_I[\bm{g}_I(\bm{w})]
  =\nabla F_S(\bm{w}).
\]
若按概率$p_i>0$非均匀抽取，则应改用
\[
  \bm{g}_I(\bm{w})
  =\frac{\nabla\ell_I(\bm{w})}{np_I}+\lambda\bm{w}.
\]
正则项梯度$\lambda\bm{w}$是确定量，无需抽样；小批量只需平均若干经过相应加权的
损失梯度，再加上这项确定梯度。小批量降低方差，却增加每步计算。

在概率工具建立前，可以先分析带确定误差的更新
\[
  \bm{w}^{(t+1)}
  =\bm{w}^{(t)}-\eta_t
  \bigl(\nabla F(\bm{w}^{(t)})+\bm{e}^{(t)}\bigr).
\]
若$F$为$L$光滑且$0<\eta_t\leq1/L$，把这一步代入下降引理，并对交叉项使用
Cauchy--Schwarz与$2ab\leq a^2+b^2$，可得
\[
  F(\bm{w}^{(t+1)})
  \leq F(\bm{w}^{(t)})
  -\frac{\eta_t}{2}\lVert\nabla F(\bm{w}^{(t)})\rVert_2^2
  +\frac{\eta_t}{2}\lVert\bm{e}^{(t)}\rVert_2^2.
\]
误差项因而必须与所求结论一起累计。固定非零偏差通常只允许迭代接近一个误差邻域；
$\sum_t\eta_t\lVert\bm{e}^{(t)}\rVert_2^2<\infty$使上述目标下降不等式中的累计
误差项有限。仅令
$\eta_t\to0$仍不够，还要核对步长总量是否足以持续纠偏，以及目标与迭代是否保持稳定。

若理想更新映射$T_t(\bm{w})=\bm{w}-\eta_t\nabla F(\bm{w})$是
$\rho_t$-Lipschitz，实际轨迹与理想轨迹之差$\delta_t$满足
\begin{equation}
  \delta_{t+1}\leq\rho_t\delta_t+\eta_t\lVert\bm{e}^{(t)}\rVert_2.
  \label{eq:opt-inexact-gradient-recurrence}
\end{equation}
展开递推可见，$\rho_t<1$时旧误差会衰减；若只有$\rho_t\leq1$，直接求和控制轨迹偏差
需要考察$\sum_t\eta_t\lVert\bm{e}^{(t)}\rVert_2$，前面的平方误差条件一般不能推出
这个一次范数条件。若$\rho_t>1$，传播因子还可能放大早期扰动。减小步长降低单步误差
影响，也会减慢理想更新。

\subsection{成对迭代与梯度自界}

同一递推可比较只替换一个样本的数据集$S,S'$。若每个样本梯度有界，且两条轨迹在使用
相同样本时由收缩映射拉近，只在抽到被替换样本时受到额外扰动，就能控制
$\lVert\bm{w}_S^{(t)}-\bm{w}_{S'}^{(t)}\rVert_2$。具体地，设每个样本梯度
$L$-Lipschitz且范数不超过$G$，两条迭代使用相同索引$I_t$。记距离为$\Delta_t$；
若$I_t$不是被替换的索引，则
\[
  \Delta_{t+1}\leq(1+\eta_tL)\Delta_t;
\]
若恰好抽到被替换项，则再增加至多$2\eta_tG$。在凸光滑且步长适当时，共同样本的梯度
更新可进一步证明为非扩张映射，从而去掉前式中的膨胀因子。这称为
\term{算法稳定性}（algorithmic stability）：训练集改动一个样本时，输出不会变化太大。
完整的期望界和随机收敛需使用第~\ref{chap:stochastic-processes-and-approximation}章的
条件期望与噪声工具。

缩放损失会等比例缩放梯度，因此学习率不能脱离损失尺度比较。梯度裁剪把过大更新限制在
指定半径内，改变了估计量并可能引入偏差。对更新式直接减去
$\eta\lambda\bm{w}$的权重衰减，在普通梯度下降中等价于$\ell_2$惩罚；使用坐标自适应缩放
时，两者通常不再等价。

\begin{lemma}[非负光滑函数的梯度自界]
\label{lem:opt-self-bounding}
若$\ell:\mathbb{R}^d\to[0,\infty)$可微且$L$光滑，则
\[
  \lVert\nabla\ell(\bm{w})\rVert_2^2\leq2L\ell(\bm{w}).
\]
\end{lemma}

\begin{proof}
在下降引理中取
$\bm{v}=\bm{w}-\nabla\ell(\bm{w})/L$，得到
\[
  \ell(\bm{v})
  \leq\ell(\bm{w})-\frac{1}{2L}
  \lVert\nabla\ell(\bm{w})\rVert_2^2.
\]
左侧非负，移项即得结论。若损失可取负值，这一步失效；给函数加常数虽不改变梯度，却会
改变右侧，所以非负性是实质条件。
\end{proof}

\section{分块、交替与分布优化}
\label{sec:opt-block-alternating-distribution}

\subsection{分块更新与约束分解}

若变量自然分成$\bm{u}$与$\bm{v}$两组，可以固定一组更新另一组。对
\[
  \Phi(\bm{u},\bm{v})
  =\frac12\lVert\bm{X}\bm{u}+\bm{Z}\bm{v}-\bm{y}\rVert_2^2
  +r(\bm{u})+s(\bm{v}),
\]
\term{分块坐标法}（block coordinate method）每次只对一个块走梯度或近端步；
\term{交替最小化}（alternating minimization）则精确或近似求解每个块的条件子问题。
若每次子问题确实降低同一个
$\Phi$，目标值单调不增。精确交替最小化时，这一点直接来自
\[
  \Phi(\bm{u}^{(t+1)},\bm{v}^{(t)})
  \leq\Phi(\bm{u}^{(t)},\bm{v}^{(t)}),\qquad
  \Phi(\bm{u}^{(t+1)},\bm{v}^{(t+1)})
  \leq\Phi(\bm{u}^{(t+1)},\bm{v}^{(t)}).
\]
但非凸情形下仍可能停在局部解或依赖初始化。

一个可复核的耦合二次例子是
\[
  \Phi(u,v)=\frac12(u-v)^2+\frac12(u-1)^2.
\]
固定$v$时的唯一最优点为$u=(v+1)/2$，固定$u$时的唯一最优点为$v=u$。从
$v^{(0)}$出发，一轮交替更新得到
$u^{(t+1)}=(v^{(t)}+1)/2$、$v^{(t+1)}=u^{(t+1)}$，因而
$v^{(t+1)}-1=(v^{(t)}-1)/2$，最终到达联合最优点$(1,1)$。这里的收敛来自这个
正定二次目标的特殊结构。换成非凸目标$(uv-1)^2$并从$(0,0)$开始时，逐块目标在一块
固定为零时没有提供离开零点的理由；若平局规则继续选择零，目标保持为$1$，而全局最优值
为$0$。所以逐块精确和目标单调都不单独保证全局最优。

约束耦合时，可以复制共享变量，再通过一致性约束分解为局部子问题。乘子负责协调副本，
这改变了每步计算位置，却没有删除原约束。若局部问题没有被足够准确地求解，或协调步长
不满足条件，分解形式本身不保证收敛。

\subsection{代理函数与单调改善}

\term{主化--最小化方法}（majorization--minimization, MM）在当前点
$\bm{w}^{(t)}$构造代理函数$Q(\bm{w}\mid\bm{w}^{(t)})$，满足
\[
  F(\bm{w})\leq Q(\bm{w}\mid\bm{w}^{(t)}),
  \qquad
  F(\bm{w}^{(t)})=Q(\bm{w}^{(t)}\mid\bm{w}^{(t)}).
\]
若$\bm{w}^{(t+1)}$使代理函数不增，则
\begin{align*}
  F(\bm{w}^{(t+1)})
  &\leq Q(\bm{w}^{(t+1)}\mid\bm{w}^{(t)})\\
  &\leq Q(\bm{w}^{(t)}\mid\bm{w}^{(t)})
  =F(\bm{w}^{(t)}).
\end{align*}
第一步用全局上界，第二步用代理改善，最后一步用切触条件。若代理只在当前点附近控制原
目标，而新点走出该邻域，单调链条便会断裂。

下降引理本身给出一个显式代理。若$F$为$L$光滑，可取
\[
  Q(\bm{w}\mid\bm{w}^{(t)})
  =F(\bm{w}^{(t)})
  +\langle\nabla F(\bm{w}^{(t)}),\bm{w}-\bm{w}^{(t)}\rangle
  +\frac{L}{2}\lVert\bm{w}-\bm{w}^{(t)}\rVert_2^2.
\]
它在当前点切触并全局上界$F$；精确最小化$Q$恰好得到步长$1/L$的梯度更新。仅让代理
与原目标在当前点具有相同梯度，却不满足全局上界，不能推出上面的单调不等式链。

变量也可以是一个分布$q$。此时“更新分布”仍是在某个可行分布集合上降低同一目标，而非
一种脱离优化的概率操作。在有限支持$\{1,\ldots,m\}$上，分布变量就是概率单纯形
\[
  \bm{\Delta}_m
  =\{\bm{q}\in\mathbb{R}^m:q_i\geq0,\ \sum_{i=1}^m q_i=1\}.
\]
归一化、非负性和支持集都是可行约束，更新后必须仍留在这个单纯形中。EM与变分方法还需要
期望、条件分布和散度。第~\ref{chap:probability-theory}章、
第~\ref{chap:mathematical-statistics}章和
第~\ref{chap:information-theory-and-statistical-geometry}章建立这些对象后，将复用
交替优化与MM的单调性核对方式。

\section{多目标与鞍点}
\label{sec:opt-multiobjective-saddle}

\subsection{支配、Pareto与共同下降}

两个任务可能产生损失向量
$\bm{f}(\bm{w})=(f_1(\bm{w}),f_2(\bm{w}))$。若一个方案两项都不差且至少一项更好，
就称它\term{支配}（dominate）另一个方案。不可被任何可行方案支配的点称为
\term{Pareto最优}（Pareto optimal）。这是一种偏序，而不是单一排名。

加权和$\alpha f_1+(1-\alpha)f_2$把偏好写成一个标量目标，$\alpha\in[0,1]$。
它能找出凸目标像边界上的支撑点，但非凸Pareto前沿的某些点无法由任何固定权重得到。
例如三个可行方案的损失向量分别为$(0,1)$、$(1,0)$和$(0.6,0.6)$。第三个方案不被
前两个方案支配，所以是Pareto最优；但对任意$\alpha\in[0,1]$，前两个方案的加权损失
中至少一个不超过$0.5$，严格小于第三个方案的$0.6$，故权重和永远不会选择第三个方案。
在当前点，方向$\bm{d}$若满足
\[
  \langle\nabla f_1(\bm{w}),\bm{d}\rangle<0,\qquad
  \langle\nabla f_2(\bm{w}),\bm{d}\rangle<0,
\]
就是\term{共同下降方向}（common descent direction）。寻找最小范数凸组合
$\alpha\nabla f_1+(1-\alpha)\nabla f_2$相当于把原点投影到两梯度的线段上。

\begin{theorem}[无共同下降方向的一阶判定]
\label{thm:opt-common-descent}
对可微的$m$个目标，在无约束点$\bm{w}$处不存在共同严格下降方向，当且仅当存在
$\alpha_i\geq0$、$\sum_i\alpha_i=1$，使
$\sum_i\alpha_i\nabla f_i(\bm{w})=\bm{0}$。
\end{theorem}

\begin{proof}
若这样的凸组合存在，而$\bm{d}$使每个内积都为负，则加权和内积也为负，与零向量矛盾。
反之，若原点不在梯度集合的凸包中，定理~\ref{thm:opt-projection-separation}给出
$\bm{d}$，使它与每个梯度的内积具有同一严格符号；必要时取$-\bm{d}$，便得到共同严格
下降方向。
\end{proof}

这个条件只排除局部的一阶共同改善，不保证全局Pareto最优；高阶变化或远处方案仍可能
同时改善所有目标。

\subsection{极小极大与鞍点间隙}

另一类问题让两个变量方向相反：
\[
  \min_{\bm{w}\in W}\max_{\bm{z}\in Z}\phi(\bm{w},\bm{z}).
\]
点$(\bm{w}^{\star},\bm{z}^{\star})$若满足
\[
  \phi(\bm{w}^{\star},\bm{z})
  \leq\phi(\bm{w}^{\star},\bm{z}^{\star})
  \leq\phi(\bm{w},\bm{z}^{\star})
\]
就称为\term{极小极大鞍点}（minimax saddle point）。这里“鞍点”描述两个变量承担
相反优化方向时的全局比较关系，不是第~\ref{sec:analysis-second-order-curvature}节中
由单个标量函数的驻点和Hessian矩阵正负曲率定义的微分鞍点。微分鞍点未必求解某个极小
极大问题；极小极大鞍点位于边界或目标不可微时，也未必满足驻点与Hessian判据。
对任意候选对定义
\[
  P(\bm{w})=\sup_{\bm{z}\in Z}\phi(\bm{w},\bm{z}),\qquad
  D(\bm{z})=\inf_{\bm{w}\in W}\phi(\bm{w},\bm{z}).
\]
总有$D(\bm{z})\leq P(\bm{w})$，所以
$P(\bm{w})-D(\bm{z})\geq0$是\term{原始--对偶间隙}
（primal--dual gap）；间隙为零时，该候选对满足鞍点不等式。

本章采用以下有限维极小极大版本：若$W,Z$非空、紧且凸，$\phi$连续，对固定
$\bm{z}$的$\phi(\cdot,\bm{z})$为凸函数，对固定$\bm{w}$的
$\phi(\bm{w},\cdot)$为凹函数，则
\[
  \min_{\bm{w}\in W}\max_{\bm{z}\in Z}\phi(\bm{w},\bm{z})
  =
  \max_{\bm{z}\in Z}\min_{\bm{w}\in W}\phi(\bm{w},\bm{z}),
\]
且两侧最值达到。紧致性提供最值的达到，凸凹性允许用分离比较混合候选，连续性保证极限
不丢失函数值。鲁棒优化中的$\bm{z}$表示最坏扰动，对抗训练让扰动方提高损失，多任务协调
则可能在权重与模型间形成极小极大结构。若缺少上述条件，两个优化次序可能不同；即使值相等，
具体梯度下降--上升仍可能旋转或发散，值的存在定理不等于算法收敛定理。

\section{双层目标与隐式更新}
\label{sec:opt-bilevel}

\subsection{展开有限次更新}

有些决策必须先经过一次学习过程才能评价。令外层变量为$\bm{\theta}$。如果内层问题有
多个最小点，只写
$\bm{w}^{\star}(\bm{\theta})\in
\argmin_{\bm{w}}L_{\mathrm{in}}(\bm{w},\bm{\theta})$
并不能定义一个单值映射，不同最小点还可能给出不同的外层损失。此时可以分别把最小点集合
上的最小、最大外层损失定义为乐观、悲观值，也可以给定由初始化和算法路径决定的单值
选择规则。下文假设内层最小解唯一，或已经给定单值选择$S$，并记
\[
  \bm{w}^{\star}(\bm{\theta})=S(\bm{\theta})
  \in\argmin_{\bm{w}}L_{\mathrm{in}}(\bm{w},\bm{\theta}).
\]
于是外层目标
\begin{equation}
  J(\bm{\theta})
  =L_{\mathrm{out}}(\bm{w}^{\star}(\bm{\theta}),\bm{\theta}).
  \label{eq:opt-bilevel-objective}
\end{equation}
这就是\term{双层优化}（bilevel optimization）。元学习把内层看成任务适应，超参数
选择让$\bm{\theta}$控制正则强度或数据权重，定向知识修正则让外层评价修改后的特定行为。

一种求导方式是把$T$步内层更新全部展开。例如
$\bm{w}^{(t+1)}=\bm{w}^{(t)}-\eta\nabla_{\bm{w}}
L_{\mathrm{in}}(\bm{w}^{(t)},\bm{\theta})$，再按链式法则从$t=T$反向传播。
展开法忠实求导于“运行$T$步”这个实际算法，但存储随$T$增长，且所得梯度对应近似内层解，
不一定对应精确最优解。

两步展开已经能看清链式结构。记
\[
  \bm{A}_t=\bm{I}-\eta\nabla_{\bm{w}\bm{w}}^2
  L_{\mathrm{in}}(\bm{w}^{(t)},\bm{\theta}),\qquad
  \bm{B}_t=\nabla_{\bm{w}\bm{\theta}}^2
  L_{\mathrm{in}}(\bm{w}^{(t)},\bm{\theta}).
\]
则
\[
  \frac{d\bm{w}^{(t+1)}}{d\bm{\theta}}
  =\bm{A}_t\frac{d\bm{w}^{(t)}}{d\bm{\theta}}-\eta\bm{B}_t,
\]
从而
\[
  \frac{d\bm{w}^{(2)}}{d\bm{\theta}}
  =\bm{A}_1\bm{A}_0\frac{d\bm{w}^{(0)}}{d\bm{\theta}}
  -\eta\bm{A}_1\bm{B}_0-\eta\bm{B}_1.
\]
外层梯度再把这个Jacobian与
$\nabla_{\bm{w}}L_{\mathrm{out}}$相乘，并加上外层目标对$\bm{\theta}$的直接偏导。
截断某一步会删除相应乘积项，因而改变所求导的算法，而不只是节省存储。

\subsection{对局部内层解隐式求导}

隐式求导跟踪的是一个选定的非退化局部最小点分支，不自动证明该分支实现全局最小值。
只有当这个分支也是全局最小点并符合前述选择规则时，它才与
$\bm{w}^{\star}(\bm{\theta})$及$J$对应。为区分两者，记局部分支及其外层目标为
$\widehat{\bm{w}}(\bm{\theta})$和
$\widehat{J}(\bm{\theta})
=L_{\mathrm{out}}(\widehat{\bm{w}}(\bm{\theta}),\bm{\theta})$。
若$L_{\mathrm{in}}$二阶连续可微，并且
\[
  \nabla_{\bm{w}}L_{\mathrm{in}}
  (\widehat{\bm{w}}(\bm{\theta}),\bm{\theta})=\bm{0},
\]
记在该分支上计算的二阶导数为
$\bm{H}_{ww}=\nabla_{\bm{w}\bm{w}}^2L_{\mathrm{in}}$和
$\bm{H}_{w\theta}=\nabla_{\bm{w}\bm{\theta}}^2L_{\mathrm{in}}$。
对驻点方程关于$\bm{\theta}$求导得到
\[
  \bm{H}_{ww}\frac{d\widehat{\bm{w}}}{d\bm{\theta}}
  +\bm{H}_{w\theta}=\bm{0}.
\]
当$\bm{H}_{ww}$可逆时，隐函数定理给出局部可微分支；对非退化局部最小点，
$\bm{H}_{ww}\succ\bm{0}$保证这一可逆性。链式法则于是给出
\begin{equation}
  \nabla_{\bm{\theta}}\widehat{J}
  =
  \nabla_{\bm{\theta}}L_{\mathrm{out}}
  -\bm{H}_{w\theta}^{\top}
  \bm{H}_{ww}^{-\top}
  \nabla_{\bm{w}}L_{\mathrm{out}}.
  \label{eq:opt-implicit-hypergradient}
\end{equation}
式~\eqref{eq:opt-implicit-hypergradient}是\term{隐式超梯度}
（implicit hypergradient）。实际计算通常解线性方程
$\bm{H}_{ww}^{\top}\bm{v}=\nabla_{\bm{w}}L_{\mathrm{out}}$，无须显式求逆。

用二次内层目标可以核对公式。设$\bm{H}\succ\bm{0}$，
\[
  L_{\mathrm{in}}(\bm{w},\bm{\theta})
  =\frac12(\bm{w}-\bm{B}\bm{\theta})^{\top}
  \bm{H}(\bm{w}-\bm{B}\bm{\theta}),\qquad
  L_{\mathrm{out}}(\bm{w})
  =\frac12\lVert\bm{C}\bm{w}-\bm{d}\rVert_2^2.
\]
这里内层最小点唯一，因此局部分支与全局选择一致：
$\widehat{\bm{w}}=\bm{w}^{\star}=\bm{B}\bm{\theta}$且
$\widehat{J}=J$。直接消元给出
\[
  \nabla_{\bm{\theta}}J
  =\bm{B}^{\top}\bm{C}^{\top}
  (\bm{C}\bm{B}\bm{\theta}-\bm{d}).
\]
隐式公式中$\bm{H}_{ww}=\bm{H}$、
$\bm{H}_{w\theta}=-\bm{H}\bm{B}$，代入后得到同一结果。计算时若线性方程残差为
$\bm{r}$，解误差为$\bm{H}^{-\top}\bm{r}$；小特征值会放大该误差。因而内层驻点残差、
线性求解残差和曲率病态是三种不同的误差来源。

若内层只近似收敛，驻点残差会进入超梯度误差；若Hessian矩阵奇异，解映射可能不唯一或不可微；
若内层非凸，所选局部分支还依赖初始化与算法路径，并且可能不是全局最小点。后一种情况下，
式~\eqref{eq:opt-implicit-hypergradient}给出的是$\widehat{J}$的梯度，而不是前述全局
双层目标$J$的梯度。展开深度、线性方程精度、内层计算预算与解选择规则因此都是问题定义的
一部分，而不只是实现细节。

\section{本章小结}
\label{sec:opt-summary}

最优化需要分别回答四个问题：候选解是否存在，某个点为何最优，算法怎样到达它，以及误差
按什么尺度收敛。连续性配合紧致性保证存在；在闭域上，下半连续性与强制性也能通过紧
下水平集保证存在。凸性把局部条件提升为全局条件，强凸、光滑性与PL条件则分别控制参数
距离、单步上界和目标误差。它们不能只用“曲率好”一语合并。

带资源约束后，KKT条件描述目标梯度与约束法向量的平衡，对偶变量提供界和敏感性解释。
Farkas证书与线性规划强对偶说明最优值相等需要证明，不能从凸性直接跳得。松弛、对偶界与
可行解共同形成最优性证书；证书说明结果离最优值多远，却不自动给出寻找结果的过程。

非光滑正则需要次梯度、投影或邻近算子，不完整梯度需要核对误差如何累计。分块、分布变量、
多目标、鞍点与双层目标改变的是变量组织和最优性概念，不能沿用同一套无条件保证。随机更新
的严格收敛将在第~\ref{chap:stochastic-processes-and-approximation}章建立条件期望与
噪声工具后补全；后续决策章节还会复用本章的对偶与鞍点语言。紧邻的
第~\ref{chap:functional-analysis-and-approximation}章将把有限维参数推广到函数空间，
区分表示存在性、逼近误差与函数类复杂度。
